%======================================================================
\chapter{Grids, the Schoenflies theorem and Moore's characterization of the plane}\label{ch:plane}
%======================================================================

We now finish the proof that a space satisfying Moore's Axioms~1--8 is homeomorphic to $\R^2$. The route is
the one recommended by Moore's own construction, organized around the Schoenflies theorem:
\begin{enumerate}
\item a compact set carrying two families of ``grid lines'' of vanishing mesh is a closed square
(Section~\ref{sec:grid}; pure topology);
\item nested double rulings (Theorem~D) give such grid lines in the closure of the interior of any quadrilateral, so
this closure is a square with the sides of the quadrilateral going to the sides of the square (the
\emph{Schoenflies theorem} in Moore's space, Section~\ref{sec:schoenflies});
\item the closed region between two nested closed curves is a collar (the \emph{annulus lemma},
Section~\ref{sec:annulus});
\item the space is exhausted by nested closed curves (Moore's Theorem~F), and the square of the first curve and the
collars between consecutive ones are pasted onto the squares $[-(n+1),n+1]^2$ (Section~\ref{sec:paste}).
\end{enumerate}
Moore's 1919 paper reaches the plane through his Theorems~E, F and~G (a ``net'' of two families of open curves);
here Theorems~E and~G are obtained as corollaries of the main theorem.

\section{A compact set with a grid is a square}\label{sec:grid}

Let $K$ be a compact subset of a metric space.

\begin{definition}\label{def:family-lines}
A \emph{family of lines} in $K$ consists of indices $i$ with positions $t_i\in[0,1]$, including an index $i_0$ with
$t_{i_0}=0$ and an index $i_1$ with $t_{i_1}=1$, and for each $i$ two closed sets $L_i$, $R_i$ (the two closed
``sides'' of the line) such that:
\begin{enumerate}[label=(\arabic*)]
\item $L_i\cup R_i=K$; $L_i=K$ if $t_i=1$ and $R_i=K$ if $t_i=0$;
\item if $t_i<t_j$ then $L_i\subseteq L_j$, $R_j\subseteq R_i$ and $L_i\cap R_j=\emptyset$;
\item the \emph{line} $\ell_i=L_i\cap R_i$ is connected, and lies in $\cl{K\setminus L_i}$ if $t_i<1$ and in $\cl{K\setminus R_i}$ if
$t_i>0$.
\end{enumerate}
A \emph{pre-grid} is a pair of families of lines $(\mathcal F_X,\mathcal F_Y)$ in $K$ such that each line of $\mathcal F_X$ meets
each line of $\mathcal F_Y$ in exactly one point. Given finitely many lines of each family, including the extreme ones, the
\emph{box} between consecutive lines $j,j'$ of $\mathcal F_X$ and $m,m'$ of $\mathcal F_Y$ is $K\setminus(L_j\cup R_{j'}\cup
L_m\cup R_{m'})$. A \emph{grid} is a pre-grid such that
\begin{enumerate}[label=(\arabic*),resume]
\item \emph{(mesh)} for every $\varepsilon>0$ and any lines $i\in\mathcal F_X$, $k\in\mathcal F_Y$ there are finite subfamilies
containing $i$, $k$ and the extreme lines all of whose boxes have closures of diameter $\le\varepsilon$;
\item \emph{(parametrized lines)} some line $k_\varphi$ of $\mathcal F_Y$ carries an injective map $\varphi\colon[0,1]\to\ell_{k_\varphi}$ with
$\varphi(s)\in L_i\iff s\le t_i$ and $\varphi(s)\in R_i\iff t_i\le s$ for all lines $i$ of $\mathcal F_X$; likewise some line of $\mathcal F_X$
carries $\psi$ compatible with the positions of $\mathcal F_Y$.
\end{enumerate}
\end{definition}

\formal{\lean{HPC.Sq.Family}, \lean{HPC.Sq.PreGrid}, \lean{HPC.Sq.box}, \lean{HPC.Sq.Grid}.}

The model: $K=[0,1]^2$, the lines of $\mathcal F_X$ the horizontal segments $y=t$ with $L=\set{y\le t}$, $R=\set{y\ge t}$, and
those of $\mathcal F_Y$ the vertical ones.

\begin{theorem}[A compact grid is a square]\label{thm:grid-square}
If the compact set $K$ carries a grid, there is a homeomorphism $f$ of $[0,1]^2$ onto $K$ sending the sides $x=0$, $x=1$,
$y=0$, $y=1$ onto the first and last lines of $\mathcal F_Y$ and of $\mathcal F_X$. Moreover $f$ can be chosen to restrict to
the parametrizations $\varphi$, $\psi$ on the corresponding lines.
\end{theorem}

\begin{proof}
For $z\in K$ define its \emph{position} in $\mathcal F_X$, $\mathrm{pos}_X(z)=\inf\set{t_i:z\in L_i}$, and similarly $\mathrm{pos}_Y$.
Then $t_i\le\mathrm{pos}_X(z)\le t_j$ whenever $z\in R_i\cap L_j$, and $\mathrm{pos}_X=t_i$ on the line $\ell_i$.

\emph{Positions of lines are dense in $[0,1]$.} The parametrized line gives $\mathrm{pos}_X(\varphi(s))=s$. If no line had its
position in an interval $(a,b)\subseteq[0,1]$, the points $\varphi(a)\ne\varphi(b)$ would lie, for every finite subfamily, in the
closure of one box (between the consecutive lines around $(a,b)$), against the mesh condition with $\varepsilon<d(\varphi(a),
\varphi(b))$.

\emph{Limit lemma.} If $\mathrm{pos}_X(z)=t_i$ then $z\in\ell_i$: by density there are lines at positions approaching $t_i$ from
both sides, and $z$ lies on their $R$-sides from below and $L$-sides from above; by (2), (3) and closedness, $z\in L_i\cap R_i$.

\emph{Continuity.} $\set{z:\mathrm{pos}_X(z)<c}=\bigcup_{t_i<c}(K\setminus R_i)$ and $\set{z:\mathrm{pos}_X(z)>c}=\bigcup_{t_i>c}
(K\setminus L_i)$, by density and the limit lemma; both are open in $K$.

\emph{Injectivity of $\Phi=(\mathrm{pos}_Y,\mathrm{pos}_X)$.} If $z,w$ have the same positions, then for every finite subfamily they lie in
the closure of the same box (or on its boundary lines), so $d(z,w)\le\varepsilon$ for every $\varepsilon$ by the mesh condition.

\emph{Surjectivity.} The crossing point of the lines $i\in\mathcal F_X$, $k\in\mathcal F_Y$ has image $(t_k,t_i)$. These points are dense in
$[0,1]^2$ and $\Phi(K)$ is compact, so $\Phi(K)=[0,1]^2$.

So $\Phi\colon K\to[0,1]^2$ is a continuous bijection from a compact space onto a Hausdorff space, a homeomorphism; let $f=\Phi^{-1}$. The
sides go where stated since $\mathrm{pos}_X(z)=0$ iff $z\in L_{i_0}$, and so on. On the parametrized line $\ell_{k_\varphi}$ we have
$\Phi(\varphi(s))=(t_{k_\varphi},s)$, which gives the last claim.
\end{proof}

\formal{\lean{HPC.Sq.Family.pos}, \lean{HPC.Sq.Grid.dense_t}, \lean{HPC.Sq.Grid.mem_L_of_pos_eq},
\lean{HPC.Sq.Grid.continuousOn_pos}, \lean{HPC.Sq.Grid.eq_of_pos_eq}, \lean{HPC.Sq.Grid.image_pos},
\lean{HPC.Sq.Grid.exists_square}, \lean{HPC.Sq.Grid.exists_square_param}.}

\begin{lemma}[Boxes are components]\label{lem:boxes}
Given finitely many lines of each family of a pre-grid (including the extreme ones), the open boxes cut out by consecutive lines
are pairwise disjoint, nonempty, and cover $K$ minus the lines. If $U\subseteq K$ is a set such that $U$ minus the lines has
exactly as many components as there are boxes, every box lies in a single component.
\end{lemma}

\begin{proof}
Disjointness and covering follow from (1)--(2); nonemptiness from the crossing of lines and (3). Each component of $U$ minus the
lines is a union of (connected pieces of) boxes; each box meets some component, and a counting argument shows that if some box met
two components there would be more components than boxes.
\end{proof}

\formal{\lean{HPC.Sq.PreGrid.box_nonempty}, \lean{HPC.Sq.PreGrid.box_sub_component}.}

\section{The Schoenflies theorem in Moore's space}\label{sec:schoenflies}

We return to a space $S$ satisfying Axioms~1--8, with a fixed metric (Proposition~\ref{prop:bridge}).

\begin{lemma}[Nested rulings]\label{lem:nested-rulings}
For any sequence $\mathcal G_1,\mathcal G_2,\dots$ of covers of $S$ by regions there are double rulings $\beta_1,\beta_2,\dots$ of a
quadrilateral $ABCDA$ such that every subdivision of $\beta_k$ lies in a region of $\mathcal G_k$, every arc of $\beta_k$ is an arc of
$\beta_{k+1}$, and every subdivision of $\beta_{k+1}$ lies in a subdivision of $\beta_k$.
\end{lemma}

\begin{proof}
$\beta_1$ is given by Theorem~D for regions (Corollary~\ref{cor:theoremD-regions}). Given $\beta_k$, every subdivision of $\beta_k$,
made into a quadrilateral in any way, is rulable with subdivisions in regions of $\mathcal G_{k+1}$ (Corollary~\ref{cor:theoremD-regions}
again); the refinement proposition (Proposition~\ref{prop:refine}) gives $\beta_{k+1}$.
\end{proof}

\formal{\lean{HPC.Moore.nested_rulings}, \lean{HPC.Moore.nested_rulings_fund}.}

\begin{theorem}[Schoenflies; Moore 1919, Theorem~E]\label{thm:schoenflies}
For every quadrilateral $ABCDA$ there is a homeomorphism $f$ of $[0,1]^2$ onto $J\cup I$ with $f(0,0)=A$, $f(0,1)=B$, $f(1,1)=C$,
$f(1,0)=D$, sending the sides $x=0$, $y=1$, $x=1$, $y=0$ of the square onto $AB$, $BC$, $CD$, $DA$. The homeomorphism can be chosen
to agree with any given parametrizations of the four sides.
\end{theorem}

\begin{proof}
Let $\mathcal G_k$ be the cover by regions of diameter $<1/k$, and $\beta_k$ the nested rulings of Lemma~\ref{lem:nested-rulings}. Let
$\mathcal A_X$ be the set of all arcs from $AB$ to $CD$ of all the $\beta_k$, and $\mathcal A_Y$ the set of all arcs from $BC$ to $DA$. Arcs of
$\mathcal A_X$ are pairwise disjoint and each meets each arc of $\mathcal A_Y$ in exactly one point (two arcs occur together in some $\beta_k$,
since the rulings are nested).

\emph{Lines.} Fix a parametrization $f_{AB}$ of $AB$ from $A$ to $B$. Record an arc $x\in\mathcal A_X$ from $X\in AB$ to $X'\in CD$ by its
\emph{line record}: the parameter $t$ of $X$, and the closures $L,R$ of the two halves $AXX'DA$ and $XBCX'X$ of the cut along $x$. Record
the side $DA$ as $(0,DA,J\cup I)$ and $BC$ as $(1,J\cup I,BC)$. These form a family of lines in $J\cup I$ (Definition~\ref{def:family-lines}):
the two sides of a cut are closed, cover $J\cup I$ and meet in the arc, which lies in the closure of each complement; along $AB$ the point
$f_{AB}(s)$ is on the $L$-side iff $s\le t$; and two disjoint arcs are nested in the order of their ends on $AB$. Together with the family
built from $\mathcal A_Y$ and $f_{BC}$ (with $AB$, $CD$ as extreme lines) this is a pre-grid, and the sides $AB$, $BC$ with $f_{AB},f_{BC}$
are parametrized lines.

\emph{Mesh.} Given a finite set of lines, all of them belong to some $\beta_k$ (together with the sides). The $(n+1)(m+1)$ boxes cut out by
the arcs of a ruling are nonempty and disjoint (Lemma~\ref{lem:boxes}) and there are exactly $(n+1)(m+1)$ subdivisions
(Proposition~\ref{prop:cells-count}); so each box lies in one subdivision, hence in a region of diameter $<1/k$. Adding lines only refines
boxes, so the mesh condition holds.

By Theorem~\ref{thm:grid-square} the compact set $J\cup I$ is a square with the stated sides. For prescribed parametrizations of all four
sides, compose with a homeomorphism of the square that fixes two adjacent sides pointwise and reparametrizes the other two
(\lean{HPC.Sq.exists_square_reparam}).
\end{proof}

\formal{\lean{HPC.Moore.lineFam}, \lean{HPC.Moore.lr_mono}, \lean{HPC.Moore.preGrid},
\lean{HPC.Moore.preGrid_box_sub_cell}, \lean{HPC.Moore.m19_schoenflies_proof}, \lean{HPC.Moore.m19_schoenflies_param},
\lean{HPC.Moore.m19_schoenflies_boundary_proof}, \lean{HPC.Moore.m19_schoenflies_boundary}.}

Theorem~E as Moore states it (two sets of arcs across $ABCDA$, one through each interior point in each direction, each arc of one set
meeting each arc of the other once) follows by taking the images of the horizontal and vertical segments under $f$
(\lean{HPC.Moore.m19_theoremE}).

\begin{corollary}\label{cor:schoenflies-curve}
For every closed curve $J$ there is a homeomorphism of $[0,1]^2$ onto $J\cup I(J)$ sending the boundary of the square onto $J$.
\end{corollary}

\begin{proof}
Choose four points on $J$ to make it a quadrilateral.
\end{proof}

\formal{\lean{HPC.Moore.exists_quad}, \lean{HPC.Moore.exists_square_of_interior}.}

\section{The annulus lemma}\label{sec:annulus}

\begin{theorem}[Annulus lemma]\label{thm:annulus}
Let $J,J'$ be closed curves with $J\cup I(J)\subseteq I(J')$. There is a homeomorphism $H$ of $J\times[0,1]$ onto the closed annulus
$(J'\cup I(J'))\setminus I(J)$ with $H(z,0)=z$ for $z\in J$ and $H(J\times\set1)=J'$.
\end{theorem}

\begin{proof}
\emph{Two cuts.} By the Schoenflies theorem for $J'$ (Corollary~\ref{cor:schoenflies-curve}) there are coordinates on $J'\cup I(J')$ in
which it is a square. On a horizontal line meeting the compact set $J\cup I(J)$, take its extreme points $P,Q$ on $J\cup I(J)$ and the
horizontal segments from them to the boundary: they give disjoint arcs $c_1=PP'$, $c_2=QQ'$ from $J$ to $J'$, each meeting $J\cup I(J)$
only in its first end and $J'$ only in its last end.

\emph{Two quadrilaterals.} Let $e$ be an arc from $P$ to $Q$ with open part in $I(J)$ (Theorem~\ref{thm:fa39-40}). Then $c_1\cup e\cup c_2$
is a cross-cut of $I(J')$, splitting it into the interiors $R_1,R_2$ of $c_1\cup e\cup c_2\cup b_i$, where $b_1,b_2$ are the arcs of $J'$
from $P'$ to $Q'$ (Corollary~\ref{cor:split-interior}). One arc $a_1$ of $J$ from $P$ to $Q$ lies in $\cl{R_1}$ and the other, $a_2$, in
$\cl{R_2}$; cutting $R_i$ along $a_i$ leaves the interior $A_i$ of the closed curve $K_i=c_1\cup b_i\cup c_2\cup a_i$. The two closed sets
$K_i\cup A_i$ fill the closed annulus and meet only along $c_1$ and $c_2$.

\emph{Gluing.} Each $K_i\cup A_i$ is a quadrilateral with sides $a_i$, $c_2$, $b_i$, $c_1$. Choose parametrizations of $c_1$, $c_2$ and of
the arcs $a_i$ (compatible at $P,Q$), and apply the Schoenflies theorem with prescribed sides (Theorem~\ref{thm:schoenflies}) to both
quadrilaterals so that the two maps agree on $c_1$ and $c_2$. Define $H(z,s)$, for $z\in a_i$, as the value at height $s$ above the
parameter of $z$ in the square of $K_i\cup A_i$. The two halves agree on $c_1,c_2$, so $H$ is a continuous bijection from the compact
$J\times[0,1]$ onto the annulus, hence a homeomorphism.
\end{proof}

\formal{\lean{HPC.Moore.exists_cuts}, \lean{HPC.Moore.annulus_decomp}, \lean{HPC.Moore.exists_collar},
\lean{HPC.Moore.m19_annulus}, \lean{HPC.Moore.m19_annulus_square}.}

Moore asserts that the annulus can be divided by a double ruling into nine cells, the central one bounded by the inner curve; this is not
needed (two quadrilaterals suffice) and is not proved in the formal development.

\section{Exhaustion and pasting}\label{sec:paste}

\begin{lemma}\label{lem:curve-about-compact}
$S$ is connected, and every compact set $M\subseteq S$ lies within some closed curve.
\end{lemma}

\begin{proof}
Connectedness is Theorem~\ref{thm:fa-elementary}(22). Cover $M$ by finitely many interiors of closed curves (Theorem~\ref{thm:fa35-36}(b))
and join each of them by an arc to a base point (Theorem~\ref{thm:fa15}); the union $M_1$ of these arcs and the closures of the interiors
is compact and connected. Cover $M_1$ by finitely many interiors of closed curves, each meeting $M_1$; their union is connected (each meets
the connected $M_1$, which they cover). By Theorem~\ref{thm:fa42} some closed curve has an interior containing all of them, hence $M$.
\end{proof}

\formal{\lean{HPC.Moore.preconnectedSpace}, \lean{HPC.Moore.exists_curve_about_compact}.}

\begin{theorem}[Moore 1919, Theorem~F]\label{thm:theoremF}
There are closed curves $J_0,J_1,\dots$ with $J_n\cup I(J_n)\subseteq I(J_{n+1})$ and $\bigcup_nI(J_n)=S$.
\end{theorem}

\begin{proof}
Let $K_1,K_2,\dots$ be the fundamental regions of Axiom~1 (every point lies in some $K_n$). Take $J_0$ about $\cl{K_0}$, and $J_{n+1}$ a
closed curve within which lies the compact set $J_n\cup I(J_n)\cup\cl{K_{n+1}}$ (Lemma~\ref{lem:curve-about-compact}).
\end{proof}

\formal{\lean{HPC.Moore.m19_theoremF_proof}, \lean{HPC.Moore.m19_theoremF}.}

\begin{proposition}[Pasting]\label{prop:paste}
Let $X$ be a Hausdorff space and $J_n,I_n\subseteq X$ ($n\ge0$) with $I_n$ open, $J_n\cap I_n=\emptyset$, $J_n\cup I_n\subseteq I_{n+1}$ and
$\bigcup I_n=X$. Suppose $J_0\cup I_0$ is the image of a homeomorphism of the square $[-1,1]^2$ sending its boundary onto $J_0$, and each
$J_n$ has a collar in $J_{n+1}\cup I_{n+1}$ ending on $J_{n+1}$ (as in Theorem~\ref{thm:annulus}). Then $X$ is homeomorphic to $\R^2$.
\end{proposition}

\begin{proof}
Inductively extend the embedding $\Psi_n$ of $[-(n+1),n+1]^2$ onto $D_n=J_n\cup I_n$ to $\Psi_{n+1}$ on $[-(n+2),n+2]^2$, mapping the square
annulus $n+1\le\nrm p\le n+2$ onto the collar by $\Psi_{n+1}(p)=H\big(\Psi_n((n+1)p/\nrm p),\,\nrm p-(n+1)\big)$ (radial extension). Then
$D_n\subseteq\interior D_{n+1}$, and the $\Psi_n$ agree where both are defined, so they define a bijection $\Psi\colon\R^2\to X$. It is
continuous (each point of $\R^2$ has a neighbourhood in some $[-(n+1),n+1]^2$) and open (each point of $X$ lies in the interior of some
$D_n$, on which $\Psi^{-1}=\Psi_{n+1}^{-1}$ is continuous).
\end{proof}

\formal{\lean{HPC.Sq.exhaust_homeomorph}, \lean{HPC.Sq.plane_of_collars}.}

\begin{theorem}[Moore's characterization of the plane]\label{thm:moore-plane}
A space satisfying Axioms~1--8 is homeomorphic to $\R^2$ (with the topology generated by its regions).
\end{theorem}

\begin{proof}
Take the curves of Theorem~\ref{thm:theoremF}. $J_0\cup I(J_0)$ is a square with boundary $J_0$ (Corollary~\ref{cor:schoenflies-curve},
rescaled to $[-1,1]^2$), and each annulus is a collar (Theorem~\ref{thm:annulus}). Proposition~\ref{prop:paste} applies, since $S$ is
Hausdorff.
\end{proof}

\formal{\lean{HPC.Moore.m19_plane_proof}, \lean{HPC.Moore.m19_plane}.}

\begin{corollary}[Moore 1919, Theorem~G and the final paragraph]\label{cor:theoremG}
\leavevmode
\begin{enumerate}[label=(\alph*)]
\item There are two sets $\mathcal G_1,\mathcal G_2$ of open curves (images of closed embeddings of $\R$) such that through each point there is
exactly one curve of each set, and each curve of $\mathcal G_1$ meets each curve of $\mathcal G_2$ in exactly one point.
\item There is a bijection $P\mapsto(x_P,y_P)$ of $S$ onto $\R^2$ such that $P$ is the sequential limit of $P_1,P_2,\dots$ (in Moore's sense)
if and only if $x_{P_n}\to x_P$ and $y_{P_n}\to y_P$.
\end{enumerate}
\end{corollary}

\begin{proof}
Transport the horizontal and vertical lines, and convergence, by the homeomorphism of Theorem~\ref{thm:moore-plane}; Moore's sequential limits are
the limits of the region topology (Proposition~\ref{prop:bridge}).
\end{proof}

\formal{\lean{HPC.Moore.m19_theoremG}, \lean{HPC.Moore.m19_openCurve}, \lean{HPC.Moore.m19_final}.}

\section*{Exercises}
\addcontentsline{toc}{section}{Exercises}

\begin{exercise}
Show that the mesh condition in Definition~\ref{def:family-lines} cannot be dropped: give a compact set with two families of ``lines'' satisfying
(1)--(3) and (5) that is not a square.
\end{exercise}

\begin{exercise}
In the plane, carry out the proof of Theorem~\ref{thm:annulus} for two concentric circles, with $c_1,c_2$ horizontal segments.
\end{exercise}

\begin{exercise}
Show that Proposition~\ref{prop:paste} fails without the hypothesis $D_n\subseteq\interior D_{n+1}$.
\end{exercise}
